% ============================================================
% §3.3  Synthesis and positioning
% ============================================================

\section{Synthesis: Open-World Autonomy and Dissertation Scope}
\label{sec:synthesis_positioning}
\label{sec:ow_gaps}
\label{sec:positioning}
\label{sec:open_world_landscape}
\label{sec:related_work_positioning}
\label{sec:perspectives_integration}
\label{sec:ch3_positioning_summary}

The literature shows that open-world autonomy is not synonymous with learning from non-identically distributed data or adapting after a task change.
Open-set and out-of-distribution methods emphasize detection relative to a reference distribution.
Concept-drift and non-stationary methods track changes in distributions, rewards, or dynamics.
Continual learning emphasizes accumulation and retention across experience, transfer learning exploits relations between source and target problems, and meta-learning acquires procedures for rapid adaptation across a task distribution.
Domain randomization and online adaptation enlarge the range of anticipated variation that an existing policy can tolerate, while replanning searches for an alternative sequence of known actions.
Each of these capabilities is relevant to an open-world agent, but none alone addresses the full case in which the assumptions or representation supporting the agent's behavior become inadequate.

This comparison reveals two broad regimes rather than a categorical separation among research communities.
When the relevant classes, variables, actions, and dimensions of variation have been anticipated, an agent can respond by updating beliefs, parameters, or policies within an adequate representation.
When they have not, continued action may require the agent to learn or revise a concept, relation, operator, transition model, skill, or goal.
Methods that support such revision confront a broader class of change, but they also expose a recurring trade-off: representational extension is often slower, more global, or more dependent on structure supplied in advance.

Foundation models change the scale of prior knowledge available to an agent but not this underlying problem.
Their broad semantic and behavioral priors can make additional situations assimilable without task-specific learning, and they can propose candidate structure when existing knowledge is insufficient.
Yet a situated agent still requires mechanisms for determining what failed, deciding what kind of revision is needed, testing candidate revisions, and grounding them in the consequences of action.
The gap is therefore not simply a lack of prior knowledge; it is the lack of reliable mechanisms that connect experienced mismatch to an appropriately scoped representational response.

The distinction between assimilation and accommodation provides a compact interpretation of this boundary without replacing the more specific problem definitions above.
Structured abstractions mediate both responses to novelty.
By encoding task-relevant invariances, they can expand the range of environmental variation that an agent assimilates within its existing knowledge and behavior.
By decomposing knowledge and behavior into identifiable, revisable components, they can make accommodation localizable when that knowledge or behavior is no longer sufficient.
A monolithic policy may register novelty only as degraded performance, whereas a structured agent can sometimes identify the failed operator, executor, object relation, or interaction variable and restrict revision to the component that no longer matches the world.

In this dissertation, we address a specific subset of the broader open-world problem through three complementary contributions.
First, systematic progress requires environments in which diverse agents can be developed and compared under controlled change.
Existing SAIL-ON-inspired environments and open-world benchmarks have made novelty injection measurable~\citep{senator2019science,holder2025introduction,balloch2022novgrid,goss2023polycraft,feeney2022novelcraft}.
In \cref{ch:benchmarks}, we present \NovelGridworlds{} and \NovelGym{}, configurable environments that support modular transformations of entities, actions, tasks, and dynamics, together with protocols and metrics for characterizing the impact of novelty and the subsequent course of adaptation.
The accompanying software ecosystem supports the development and evaluation of modular agent architectures designed for open-world settings.
This contribution operationalizes the study of how novelty affects existing knowledge and behavior and when an adaptive response becomes necessary.

Second, open-world agents require mechanisms for localized accommodation after execution failure.
Replanning can rearrange known operators, and undirected exploration can discover behavior in principle, but neither by itself identifies the precise missing behavior when an expected action effect no longer occurs.
In \cref{ch:rapid_learn}, we study action abstractions in which symbolic operators are associated with executable behaviors.
When an executor fails, the intended effects of its operator define what recovery must accomplish, converting global task failure into a targeted skill-learning problem while preserving unaffected knowledge and behaviors.

Third, agents should not be forced to accommodate every change they encounter.
Some apparent novelty results from representations that are unnecessarily tied to a particular object, physical condition, robot, or control space.
In \cref{ch:flex,ch:pho}, we study interaction abstractions that represent manipulation at the level of object interaction rather than robot-specific control.
By capturing forces, object dynamics, and other task-relevant invariances, these abstractions allow existing skills to generalize across objects, physical parameters, and embodiments; in the terminology of this chapter, they expand the range of variation that can be assimilated without structural revision or retraining.

The action and interaction contributions therefore instantiate complementary consequences of the framework taken from \cref{ch:novelty_open_world}.
Action abstractions provide structured interfaces for accommodation when novelty invalidates an existing behavior.
Interaction abstractions enlarge the domain of assimilation by making some environmental and embodiment variation irrelevant to the learned skill.
The benchmark contribution makes these responses available for controlled study and evaluation.
Together, they support the thesis that representational structure shapes whether environmental change produces brittle failure, targeted accommodation, or continued task performance through generalization.

These contributions address only part of the broader problem of open-world autonomy.
They do not solve general open-world perception, autonomously construct arbitrary ontologies, or integrate all four requirements into a single lifelong agent.
The main technical contributions assume access to task-relevant state and do not address safety-critical continual adaptation, where violations during learning may be unacceptable~\citep{coursey2026safe}.
They also do not address the normative conflicts that arise when agents operating among people encounter novel clashes among goals, rules, or constraints~\citep{jones2026oamncc}.
Finally, the methods are primarily failure-driven and do not treat novelty principally as an opportunity for self-directed exploration or goal formation.
Instead, they isolate and study how explicit structure can support evaluation, localized behavioral accommodation, and generalization across important classes of variation.
These boundaries define the scope of the dissertation rather than the scope of open-world autonomy as a whole.
The next chapter formalizes the agent--environment setting, the notion of novelty, and the action and interaction abstractions used to develop these claims in the technical chapters.
